\section{Conclusion and Perspectives}

In this paper, we have investigated an optimal control problem governed by a stochastic 
advection--diffusion equation driven by additive noise.

We first introduced the mathematical model and derived its variational formulation within 
an appropriate Hilbert space framework. Under suitable assumptions on the diffusion coefficient, 
the velocity field, the control variable, and the initial condition, we established the 
existence and uniqueness of weak solutions to the stochastic state equation.

The optimal control problem was then formulated through a quadratic tracking-type cost 
functional. Using standard tools from the calculus of variations and functional analysis, 
we proved the existence of optimal controls. The derivation of the adjoint equation enabled 
us to obtain first-order necessary optimality conditions and to characterize the optimal control 
through an optimality system coupling the state equation, the adjoint equation, and the control 
law.

A finite element discretization based on continuous piecewise linear elements was proposed for 
the numerical approximation of the optimality system. The resulting discrete formulation provides 
a practical framework for implementation within the FEniCSx environment. Numerical experiments 
were presented to illustrate the theoretical developments and to demonstrate the effectiveness 
of the proposed optimization strategy.

The methodology developed in this work is applicable to a wide range of stochastic transport 
phenomena arising in environmental sciences, groundwater management, atmospheric pollution 
control, and climate modeling.

Several extensions of the present work may be considered. A first direction consists in studying 
stochastic advection--diffusion equations driven by multiplicative noise. Another interesting 
perspective concerns the analysis of control constraints and sparsity-promoting cost functionals. 
From the numerical viewpoint, future work will focus on a rigorous error analysis of the finite 
element approximation and on the development of efficient parallel algorithms for large-scale 
simulations. Finally, the extension of the proposed framework to nonlinear stochastic transport 
equations and coupled multiphysics systems constitutes a promising research direction.
